\documentclass[otchet]{SCWorks}
% Тип обучения (одно из значений):
%    bachelor   - бакалавриат (по умолчанию)
%    spec       - специальность
%    master     - магистратура
% Форма обучения (одно из значений):
%    och        - очное (по умолчанию)
%    zaoch      - заочное
% Тип работы (одно из значений):
%    coursework - курсовая работа (по умолчанию)
%    referat    - реферат
%  * otchet     - универсальный отчет
%  * nirjournal - журнал НИР
%  * digital    - итоговая работа для цифровой кафедры
%    diploma    - дипломная работа
%    pract      - отчет о научно-исследовательской работе
%    autoref    - автореферат выпускной работы
%    assignment - задание на выпускную квалификационную работу
%    review     - отзыв руководителя
%    critique   - рецензия на выпускную работу

% * Добавлены вручную. За вопросами к @mchernigin
\usepackage{preamble}

\begin{document}

% Кафедра (в родительном падеже)
\chair{математической кибернетики и компьютерных наук}

% Тема работы
\title{Анализ двоичного дерева поиска}

% Курс
\course{2}

% Группа
\group{211}

% Факультет (в родительном падеже) (по умолчанию "факультета КНиИТ")
% \department{факультета КНиИТ}

% Специальность/направление код - наименование
 \napravlenie{02.03.02 "--- Фундаментальная информатика и информационные технологии}
% \napravlenie{02.03.01 "--- Математическое обеспечение и администрирование информационных систем}
% \napravlenie{09.03.01 "--- Информатика и вычислительная техника}
%\napravlenie{09.03.04 "--- Программная инженерия}
% \napravlenie{10.05.01 "--- Компьютерная безопасность}

% Для студентки. Для работы студента следующая команда не нужна.
% \studenttitle{студентки}

% Фамилия, имя, отчество в родительном падеже
\author{Бурова Арсения Александровича}

% Заведующий кафедрой 
\chtitle{доцент, к.\,ф.-м.\,н.}
\chname{С.\,В.\,Миронов}

% Руководитель ДПП ПП для цифровой кафедры (перекрывает заведующего кафедры)
% \chpretitle{
%     заведующий кафедрой математических основ информатики и олимпиадного\\
%     программирования на базе МАОУ <<Ф"=Т лицей №1>>
% }
% \chtitle{г. Саратов, к.\,ф.-м.\,н., доцент}
% \chname{Кондратова\, Ю.\,Н.}

% Научный руководитель (для реферата преподаватель проверяющий работу)
\satitle{доцент, к.\,ф.-м.\,н.} %должность, степень, звание
\saname{Ю.\,Н.\,Кондратова}

% Руководитель практики от организации (руководитель для цифровой кафедры)
\patitle{доцент, к.\,ф.-м.\,н.}
\paname{С.\,В.\,Миронов}

% Руководитель НИР
\nirtitle{доцент, к.\,п.\,н.} % степень, звание
\nirname{В.\,А.\,Векслер}

% Семестр (только для практики, для остальных типов работ не используется)
\term{2}

% Наименование практики (только для практики, для остальных типов работ не
% используется)
\practtype{учебная}

% Продолжительность практики (количество недель) (только для практики, для
% остальных типов работ не используется)
\duration{2}

% Даты начала и окончания практики (только для практики, для остальных типов
% работ не используется)
\practStart{01.07.2024}
\practFinish{13.01.2024}

% Год выполнения отчета
\date{2026}

\maketitle

% Включение нумерации рисунков, формул и таблиц по разделам (по умолчанию -
% нумерация сквозная) (допускается оба вида нумерации)
\secNumbering

\tableofcontents

% Раздел "Обозначения и сокращения". Может отсутствовать в работе
% \abbreviations
% \begin{description}
%     \item ... "--- ...
%     \item ... "--- ...
% \end{description}

% Раздел "Определения". Может отсутствовать в работе
% \definitions

% Раздел "Определения, обозначения и сокращения". Может отсутствовать в работе.
% Если присутствует, то заменяет собой разделы "Обозначения и сокращения" и
% "Определения"
% \defabbr

\section{Реализация}

\subsection{Дерево и функции для него}
\begin{Verbatim}[
  numbers=left,
  breaklines=true,
  breakanywhere=true
]

struct tree {
    int inf;
    tree* left;
    tree* right;
    tree* parent;
};

tree* node(int x) {
    tree* n = new tree;
    n->inf = x;
    n->left = n->right = n->parent = NULL;
    return n;
}

//вставка элемента
void insert(tree*& tr, int x) {
    tree* n = node(x);
    if (!tr) tr = n; //если дерево пустое - корень
    else {
        tree* y = tr;
        while (y) { //ищем куда вставлять
            if (n->inf > y->inf) { //правая ветка
                if (y->right) y = y->right;
                else {
                    n->parent = y; //узел становится правым ребенком
                    y->right = n;
                    break;
                }
            }
            else if (n->inf < y->inf) { //левая ветка
                if (y->left) y = y->left;
                else {
                    n->parent = y; //узел становится левым ребенком
                    y->left = n;
                    break;
                }
            }
            else {
                delete n;
                break;
            }
        }
    }
}

//поиск элемента в дереве бинарного поиска
tree* find(tree* tr, int x) {
    if (!tr || tr->inf == x) return tr; //нашли или дошли до конца ветки
    if (x < tr->inf)
        return find(tr->left, x); //ищем по левой ветке
    else
        return find(tr->right, x); //ищем по правой ветке
}

tree* Min(tree* tr) {
    if (!tr->left) return tr; //нет левого ребенка
    else return Min(tr->left); //идем по левой ветке до конца
}

tree* Max(tree* tr) {
    if (!tr->right) return tr; //нет правого ребенка
    else return Max(tr->right); //идем по правой ветке до конца
}


tree* Next(tree* tr, int x) {//поиск следующего
    tree* n = find(tr, x);
    if (n->right)//если есть правый ребенок
        return Min(n->right);//min по правой ветке
    tree* y = n->parent; //родитель
    while (y && n == y->right) {//пока не дошли до корня или узел - правый ребенок
        n = y;//идем вверх по дереву
        y = y->parent;
    }
    return y;//возвращаем родителя
}

void Delete(tree*& tr, tree* v) {
    tree* p = v->parent;

    if (!v->left && !v->right) {//нет детей
        if (!p) { //корень без детей
            tr = NULL;
        }
        else if (p->left == v)
            p->left = NULL;
        else
            p->right = NULL;
        delete v;
    }
    else if (!v->left || !v->right) {//один ребенок
        if (!p) { //корень с одним ребенком
            if (v->left) {
                tr = v->left;
                v->left->parent = NULL;
            }
            else {
                tr = v->right;
                v->right->parent = NULL;
            }
            delete v;
        }
        else {
            if (v->left) {
                if (p->left == v)
                    p->left = v->left;
                else
                    p->right = v->left;
                v->left->parent = p;
            }
            else {
                if (p->left == v)
                    p->left = v->right;
                else
                    p->right = v->right;
                v->right->parent = p;
            }
            delete v;
        }
    }
    else {//удаление узла с двумя детьми
        tree* s = Next(tr, v->inf); //находим следующий за удаляемым узел
        v->inf = s->inf;
        tree* parent = s->parent;
        if (parent->left == s) {
            parent->left = s->right;
            if (s->right)
                s->right->parent = parent;
        }
        else {
            parent->right = s->right;
            if (s->right)
                s->right->parent = parent;
        }
        delete s;
    }
}
\end{Verbatim}

\subsection{Обходы}

\begin{Verbatim}[
  numbers=left,
  breaklines=true,
  breakanywhere=true
]
void preorder(tree* tr) {
    if (tr) {
        cout << tr->inf << " "; //корень
        preorder(tr->left); //левое поддерево
        preorder(tr->right); //правое поддерево
    }
}
//симметричный обход (левый-корень-правый)
void inorder(tree* tr) {
    if (tr) {
        inorder(tr->left); //левое поддерево
        cout << tr->inf << " "; //корень
        inorder(tr->right); //правое поддерево
    }
}
//обратный обход (левый-правый-корень)
void postorder(tree* tr) {
    if (tr) {
        postorder(tr->left); //левое поддерево
        postorder(tr->right); //правое поддерево
        cout << tr->inf << " "; //корень
    }
}
\end{Verbatim}



\newpage

\section{Анализ}

Пусть \(n\) — количество узлов в дереве, \(h\) — высота дерева.

\subsection{Вставка элемента}

Функция \texttt{insert} выполняет итеративный спуск от корня до места вставки. На каждом шаге выполняется сравнение и переход к левому или правому поддереву. Количество итераций не превышает высоты дерева \(h\). Создание узла и установка указателей выполняются за \(O(1)\).

\[
T_{\text{insert}}(h) = O(h)
\]

\subsection*{Поиск элемента}
Рекурсивная реализация: на каждом уровне рекурсии выполняется переход в левое или правое поддерево. В худшем случае (элемент отсутствует или находится в листе) число вызовов равно \(h\).
\[
T_{\text{find}}(h) = O(h)
\]

\subsection{Поиск минимума и максимума}
Рекурсивный спуск по левой (или правой) ветви до листа. Количество шагов не превышает \(h\).
\[
T_{\min/\max}(h) = O(h)
\]

\subsection{Поиск следующего элемента}
Алгоритм:
\begin{enumerate}
    \item Поиск узла по ключу — \(O(h)\).
    \item Если есть правый потомок — поиск минимума в правом поддереве за \(O(h)\).
    \item Иначе — подъём вверх по указателям \texttt{parent}, в худшем случае до корня за \(O(h)\).
\end{enumerate}
\[
T_{\text{Next}}(h) = O(h)
\]

\subsection{Удаление элемента}
\begin{itemize}
    \item Лист: обновление указателя родителя — \(O(1)\).
    \item Один потомок: перенаправление указателей — \(O(1)\).
    \item Два потомка: поиск следующего узла (минимум в правом поддереве) — \(O(h)\), затем удаление узла с не более чем одним потомком — \(O(1)\).
\end{itemize}
В худшем случае (удаление узла с двумя потомками):
\[
T_{\text{Delete}}(h) = O(h)
\]

\subsection{Обходы дерева (preorder, inorder, postorder)}
Каждый узел посещается ровно один раз. Общее количество рекурсивных вызовов пропорционально \(n\).
\[
T(n) = O(n)
\]

\subsection{Расход памяти}

\paragraph{Память для хранения дерева:}
Каждый узел структуры \texttt{tree} содержит:
\begin{itemize}
    \item \texttt{int inf} — 4 байта;
    \item \texttt{tree* left} — указатель (4 или 8 байт);
    \item \texttt{tree* right} — указатель (4 или 8 байт);
    \item \texttt{tree* parent} — указатель (4 или 8 байт).
\end{itemize}
Итого на один узел: \(4 + 3 \times \text{sizeof(pointer)}\) байт. Для \(n\) узлов:
\[
M_{\text{storage}}(n) = O(n)
\]

\paragraph{Дополнительная память на стеке вызовов:}
Рекурсивные функции (\texttt{find}, \texttt{Min}, \texttt{Max}, обходы) используют стек вызовов глубиной до \(h\) в худшем случае. Таким образом:
\[
M_{\text{stack}}(h) = O(h)
\]

\paragraph{Общая пространственная сложность:}
\[
M(n, h) = O(n) + O(h) = O(n + h)
\]

\end{document}
